\documentclass[conference]{IEEEtran}
\usepackage{amsmath,graphicx}
\title{Vision-Guided Selective UV Curing for Gel Nails}
\author{\IEEEauthorblockN{Author names pending}}
\begin{document}
\maketitle
\section{Introduction}
Full-hand gel nail lamps illuminate both nails and adjacent skin. Selective curing restricts illumination to the vicinity of each nail. This project, developed in Arequipa, Peru, proposes modular vision-to-motion software, a low-cost arm with collimated UV illumination and planned hardware interlocks, and a quantitative comparison of illuminated skin area.
\section{Related Work}
Zhivagui et al. examined DNA damage in cultured cells after irradiation with a nail polish dryer \cite{zhivagui}. MediaPipe Hands estimates hand landmarks from RGB images \cite{mediapipe}. Distal landmarks serve as spatial priors in the contour baseline. The primary pipeline uses a published YOLOv8s instance-segmentation model specialized for nails \cite{nailseg}. Planar joint angles are calculated with a standard two-link kinematic model \cite{lynch}.
\section{System Design}
\setcounter{subsection}{2}
\subsection{Fingernail Segmentation and Homography Mapping}
The ONNX model returns candidate nail regions, detection scores, and mask coefficients. Candidate boxes are filtered by non-maximum suppression and instance masks are reconstructed from the mask prototypes. Centroids and equivalent radii are extracted from the binary masks.
\begin{equation}
s[x,y,1]^\mathsf{T}=H[u,v,1]^\mathsf{T}.
\end{equation}
Here $(u,v)$ denotes image pixels, $(x,y)$ denotes pad coordinates in millimeters, $H$ is a homography fitted from four correspondences, and $s$ is a homogeneous scale. Calibration error is evaluated on independent grid points.
\begin{align}
\cos\theta_2 &= \frac{x_r^2+y_r^2-L_1^2-L_2^2}{2L_1L_2},\\
\theta_1 &= \operatorname{atan2}(y_r,x_r)-\operatorname{atan2}(L_2\sin\theta_2,L_1+L_2\cos\theta_2).
\end{align}
Here $(x_r,y_r)$ is the target in millimeters relative to the arm origin, $L_1,L_2$ are link lengths in millimeters, $\theta_1$ is the base angle, and $\theta_2$ is the relative elbow angle. Targets are screened against workspace and joint limits before transmission.
\section{Evaluation Protocol}
Segmentation is evaluated by intersection over union against annotated masks. Targeting error is the planar distance from the measured spot center to the annotated centroid. Targeting errors and curing durations are summarized by the sample mean and standard deviation.
\begin{equation}
\mathrm{SERR}=100(1-A_{\mathrm{sel}}/A_{\mathrm{base}}).
\end{equation}
SERR denotes skin exposure reduction ratio in percent. $A_{\mathrm{base}}$ is the skin area within the conventional illumination region. $A_{\mathrm{sel}}$ is the skin area intersecting the union of the circular spots. Non-nail pixels in the hand bounding box are used as an area proxy when a hand mask is unavailable.
% Fig. 2: insert annotated capture and projected spots.
% Table II and Fig. 3: insert experimental area comparison after trials.
\section{Limitations}
The method assumes a stationary hand, planar calibration, and constant spot diameter. Height variations can introduce projection error. Contour refinement depends on image contrast. Dose characterization requires irradiance measurements.
\section{Conclusion and Future Work}
The software converts image targets into a serial curing sequence. Further work includes evaluation on samples from Arequipa, closed-loop tracking, and adaptive optics.
\begin{thebibliography}{9}
\bibitem{zhivagui} M. Zhivagui et al., DNA damage and somatic mutations in mammalian cells after irradiation with a nail polish dryer, \emph{Nature Communications}, vol. 14, art. 276, 2023, doi: 10.1038/s41467-023-35876-8.
\bibitem{nailseg} austingg, Finger nail segmentation, GitHub repository, \texttt{github.com/austingg/finger-nail-seg}.
\bibitem{mediapipe} F. Zhang et al., MediaPipe Hands: On-device real-time hand tracking, arXiv:2006.10214, 2020.
\bibitem{lynch} K. M. Lynch and F. C. Park, \emph{Modern Robotics: Mechanics, Planning, and Control}. Cambridge Univ. Press, 2017.
\end{thebibliography}
\end{document}
